\documentclass[11pt,a4paper]{article}
\usepackage[utf8]{inputenc}
\usepackage[T1]{fontenc}
\usepackage[french]{babel}
\frenchsetup{StandardItemLabels=true}
\usepackage{lmodern,microtype}
\usepackage[margin=20mm,headheight=15pt,headsep=7mm,footskip=10mm]{geometry}
\usepackage{amsmath,amssymb,booktabs,array,tabularx,graphicx,lastpage}
\usepackage{xcolor,fancyhdr,titlesec,enumitem}
\usepackage[unicode,colorlinks=true,linkcolor=blue!45!black,urlcolor=blue!45!black,citecolor=blue!45!black]{hyperref}
\definecolor{navy}{HTML}{18364B}
\hypersetup{pdftitle={Prédire le total en force athlétique}}
\pagestyle{fancy}
\fancyhf{}
\fancyhead[L]{\small\color{navy}Prédire le total en force athlétique}
\fancyhead[R]{\small H. Tchani Yacoubou, S. Thibault}
\fancyfoot[C]{\small\thepage\,/\,\pageref{LastPage}}
\renewcommand{\headrulewidth}{0.3pt}
\setlength{\parindent}{0pt}
\setlength{\parskip}{4pt}
\setlength{\emergencystretch}{2em}
\titleformat{\section}{\large\bfseries\color{navy}}{\thesection.}{0.5em}{}
\titleformat{\subsection}{\normalsize\bfseries\color{navy}}{\thesubsection}{0.5em}{}
\titlespacing*{\section}{0pt}{8pt}{3pt}
\titlespacing*{\subsection}{0pt}{5pt}{2pt}
\setlist{nosep,leftmargin=1.4em,topsep=2pt}
\setlist[itemize]{label=\textbullet}
\newcommand{\code}[1]{\texttt{\small #1}}
\newcommand{\kg}{\,kg}
\newcounter{fig}
\newcommand{\reportfig}[4][\linewidth]{\begin{center}\includegraphics[width=#1]{figures/#2}\refstepcounter{fig}\label{#3}\par{\small\textbf{Figure~\thefig.} #4\par}\end{center}\vspace{-6pt}}
\newenvironment{tableau}{\par\smallskip\noindent\begin{minipage}{\linewidth}\centering}{\end{minipage}\par\smallskip}
\newcounter{tab}
\newcommand{\legende}[2]{\refstepcounter{tab}\label{#1}{\small\textbf{Tableau~\thetab.} #2\par}}

\begin{document}

\begin{center}
{\LARGE\bfseries\color{navy}Prédire le total en force athlétique}\par
\vspace{3pt}
{\large Modèles individuels, erreurs communes et correction par appariement}\par
\vspace{5pt}
Hamdane Tchani Yacoubou et Solène Thibault\par
{\small Master 2 Mathématiques et IA, Université Paris-Saclay. Apprentissage supervisé, challenge de régression.}\par
{\small Code : \href{https://github.com/Hamdane11/performances-powerlifting}{github.com/Hamdane11/performances-powerlifting}}
\end{center}

\noindent\textbf{Résumé.} Nous prédisons le total soulevé en compétition de force athlétique à partir de
16 variables décrivant l'athlète et la compétition. Nous étudions d'abord quatre familles de modèles
(régression Ridge, boosting d'arbres, réseau TabM, modèle préentraîné TabPFN) et le choix de leurs
réglages ; le meilleur, TabPFN, atteint 76,94\kg{} de RMSE en validation croisée. Tous les modèles se
trompent sur les mêmes athlètes et dans le même sens, ce qui explique pourquoi les combiner rapporte peu
(0,3\kg). Nous proposons alors une \emph{correction par appariement} : pour chaque ligne, nous
cherchons dans les données étiquetées les participations qui ont de bonnes chances d'être celles du même
athlète, et un modèle de second niveau corrige la prédiction à partir de leurs totaux. Appliquée à notre
mélange de modèles, elle réduit l'erreur de 0,37\kg{} en validation croisée imbriquée et de 0,68\kg{} sur un
test interne de 21\,000 lignes. Notre solution finale obtient 76,27\kg{} en validation croisée, 75,83\kg{}
sur le test interne et 75,65\kg{} sur le classement public.

% =====================================================================
\section{Le problème}\label{sec:probleme}

\subsection{Données et métrique}
Chaque ligne décrit la participation d'un athlète à une compétition complète (squat, développé couché,
soulevé de terre). La cible \code{TotalKg} est la somme des meilleurs essais réussis. Nous disposons de
105\,000 lignes étiquetées et devons prédire 45\,000 lignes de test. La métrique est la racine de
l'erreur quadratique moyenne,
\[
\mathrm{RMSE} = \Big(\tfrac1n \textstyle\sum_{i=1}^n (y_i - \hat y_i)^2\Big)^{1/2},
\]
en kilogrammes. Le classement public de Kaggle porte sur environ 30\,\% du test, le classement final sur
les 70\,\% restants.

Les 16 variables explicatives sont de trois types :
\begin{itemize}
\item \textbf{l'athlète} : sexe (69\,\% d'hommes), âge (médiane 25 ans), poids de corps (médiane 81\kg),
pays et région, tranches d'âge et d'année de naissance ;
\item \textbf{la catégorie} : catégorie de poids (146 modalités, car les fédérations n'ont pas les mêmes),
équipement (\emph{Raw} 62\,\%, \emph{Single-ply} 23\,\%, \emph{Wraps} 12\,\%, \emph{Multi-ply} 2\,\%),
contrôle antidopage (76\,\% de compétitions contrôlées) ;
\item \textbf{la compétition} : fédération (399 modalités), fédération de tutelle, pays, région, année
(1966 à 2026), homologation.
\end{itemize}
Il n'y a pas de valeur manquante : l'absence d'information est codée \code{Unknown} (région de l'athlète
inconnue dans 64\,\% des lignes). Aucune variable n'identifie l'athlète.

\subsection{Ce que l'on peut prédire simplement}
Le tableau~\ref{tab:references} prédit chaque total par la moyenne de son groupe, calculée sur d'autres
lignes que celle prédite. La colonne $R^2 = 1 - \mathrm{RMSE}^2/\mathrm{Var}(y)$ donne la part de la dispersion
des totaux que la prédiction explique : prédire tout le monde par la moyenne générale laisse toute la
dispersion ($R^2 = 0$). Connaître le sexe, l'équipement et la catégorie de poids en explique déjà les deux
tiers : l'erreur passe de 171 à 100\kg, et son carré de 29\,300 à 10\,100 (figure~\ref{fig:donnees}). Le tiers
restant, ce sont les différences entre athlètes d'un même groupe ; c'est là que se joue le problème.

\begin{tableau}\small
\begin{tabular}{lcc}
\toprule
Prédiction & RMSE (kg) & $R^2$ \\
\midrule
Moyenne générale & 171,3 & 0 \\
Moyenne par sexe & 128,3 & 0,44 \\
Moyenne par sexe et catégorie de poids & 105,5 & 0,62 \\
Moyenne par sexe, équipement et catégorie de poids & 100,3 & 0,66 \\
Notre solution finale & 76,3 & 0,80 \\
\bottomrule
\end{tabular}\par\vspace{2pt}
\legende{tab:references}{Prédictions de référence, sur les mêmes lignes et les mêmes plis que nos modèles
(validation croisée sur 84\,000 lignes, section~\ref{sec:protocole}).}
\end{tableau}

\reportfig{fig_donnees}{fig:donnees}{À gauche, distribution du total. À droite, total selon le poids de
corps (6\,000 athlètes tirés au hasard) et médiane par tranche de 5\kg.}

% =====================================================================
\section{Protocole de validation}\label{sec:protocole}
Nous mettons de côté 21\,000 des 105\,000 lignes étiquetées comme \textbf{test interne}, mesuré une seule
fois à la fin (il a servi à évaluer nos toutes premières versions, mais aucun réglage de la solution finale
n'a été choisi en le regardant). Sur les 84\,000 autres lignes, chaque modèle est évalué en
\textbf{validation croisée à 5 plis} : il apprend sur 67\,200 lignes et prédit les 16\,800 autres ; tous les
modèles utilisent les mêmes plis. Les réglages sont choisis sur un découpage de ces 84\,000 lignes en
63\,000 lignes d'apprentissage et 21\,000 de validation. Ce qui est ajusté sur les prédictions hors pli
elles-mêmes (poids d'un mélange, correction de second niveau) est évalué en \textbf{validation croisée
imbriquée} : pour mesurer l'erreur sur un pli, l'ajustement est refait sans aucune ligne de ce pli. Les écarts
entre méthodes sont donnés avec un intervalle de confiance à 95\,\% (bootstrap apparié sur les lignes). Les
modèles soumis sont réentraînés sur les 105\,000 lignes.

% =====================================================================
\section{Prétraitements}\label{sec:pretraitements}

\begin{itemize}
\item \textbf{Âges.} Les 18 âges inférieurs à 5 ans sont mis en valeur manquante. Les colonnes reprennent le
format de la base publique OpenPowerlifting, où un âge en ,5 (48\,\% des lignes) signifie que seule l'année de
naissance est connue : 19,5 veut dire « 19 ou 20 ans »~[16]. Nous en tirons un indicateur et une année de
naissance estimée.
\item \textbf{Catégories.} La modalité \code{Unknown} est gardée telle quelle : elle n'est pas due au hasard
(région inconnue pour 95\,\% des athlètes hors des États-Unis, 28\,\% aux États-Unis).
\item \textbf{Variables ajoutées.} Ridge et TabPFN reçoivent les 16 variables d'origine. TabM (24 colonnes) reçoit
en plus la limite de la catégorie de poids, son caractère ouvert (\code{120+}) ou fermé, l'écart et le rapport
entre le poids de corps et cette limite, l'année de naissance estimée, l'indicateur d'âge et les combinaisons
pays\,|\,région. LightGBM (32 colonnes) et CatBoost (36) reçoivent en outre des croisements de catégories
(sexe $\times$ équipement, fédération $\times$ année, etc.) et des indicateurs de précision de la pesée
(poids au kilo, au demi-kilo ou au dixième près).
\item \textbf{Cible ramenée à un poids de référence.} Sous similitude géométrique, la force est proportionnelle
à la section des muscles, donc à la masse à la puissance $2/3$~[14]. Certains modèles apprennent $z = y/s$,
avec $s = (b/80)^p$ et $b$ le poids de corps, puis prédisent $s\,\hat z$. Avec des poids d'apprentissage $s^2$,
l'objectif reste exactement l'erreur quadratique en kilogrammes : $\sum_i s_i^2 (z_i - \hat z_i)^2 =
\sum_i (y_i - s_i \hat z_i)^2$.
\end{itemize}

% =====================================================================
\section{Les modèles individuels}\label{sec:modeles}

\subsection{Régression Ridge}
$\hat y = X\hat\beta$ avec $\hat\beta = \arg\min_\beta \|y - X\beta\|^2 + \alpha\|\beta\|^2$~[1], les catégories
codées en indicatrices et les variables numériques standardisées. La pénalité change peu le résultat
(86,1\kg{} pour $\alpha \in [1, 10]$, 86,6 pour $\alpha = 100$). La forme de la cible compte davantage. Ridge
est additif : chaque variable ajoute un nombre fixe de kilogrammes, le même quel que soit le poids de corps.
Or les écarts grandissent avec le poids : chez les hommes, les athlètes en \emph{Single-ply} soulèvent en
moyenne 44\kg{} de plus que ceux en \emph{Raw} entre 55 et 70\kg{} de poids de corps, et 90\kg{} de plus
entre 110 et 150\kg ; en proportion, l'écart reste presque constant (de +10 à +14\,\%). Avec la cible ramenée,
l'objectif devient $\sum_i \big(y_i - s_i(\beta_0 + \beta^\top x_i)\big)^2 + \alpha\|\beta\|^2$ : c'est une régression
Ridge sur les variables multipliées par $s$, où chaque effet devient proportionnel à $s$. L'erreur passe de
86,6\kg{} ($p = 0$) à 84,4 ($p = 2/3$) et 83,9 ($p = 1$). Pour LightGBM, la même transformation ne change rien
($+0{,}08$\kg, intervalle $[-0{,}02 ; 0{,}18]$) : un arbre peut couper sur le poids de corps et apprendre un
effet différent dans chaque tranche.

\subsection{Boosting d'arbres : LightGBM et CatBoost}
LightGBM et CatBoost sont des boosting de gradient~[2] avec la perte quadratique : chaque arbre est ajusté aux
résidus du modèle courant puis ajouté avec un taux $\eta < 1$ ; une pénalité $\lambda$ rapproche de zéro la
valeur des feuilles qui contiennent peu de lignes~[3], et le nombre d'arbres est fixé par early stopping.

\textbf{LightGBM}~[4] fait pousser les arbres feuille par feuille et traite les catégories directement. Pour
mesurer l'effet de ses réglages, nous avons tiré au hasard 60 configurations de huit hyperparamètres (taille
des arbres, nombre minimal de lignes par feuille, $\lambda$, fractions de lignes et de variables, seuils
aléatoires, lissage des catégories, exposant $p$), chacune entraînée sur 63\,000 lignes avec $\eta = 0{,}05$,
le nombre d'arbres étant fixé par early stopping sur les 21\,000 lignes de validation. Les 61 erreurs de
validation (60 configurations tirées et notre configuration de départ : 127 feuilles, $\lambda = 10$)
s'étalent de 77,8 à 78,9\kg, avec un écart-type de 0,26 (figure~\ref{fig:lgb}). Un seul réglage a un effet net : sous-échantillonner les
lignes pour chaque arbre dégrade l'erreur (médiane 78,54\kg{} avec 70\,\% des lignes, 78,16 sans
sous-échantillonnage). Réentraînées avec $\eta = 0{,}02$ et mesurées en validation croisée, la meilleure
configuration (15 feuilles, $\lambda = 30$, 3\,974 arbres) et la configuration de départ (1\,035 arbres)
obtiennent 78,10 et 78,11\kg, contre 78,23 pour la configuration de départ avec $\eta = 0{,}04$. Le gain
vient donc du taux d'apprentissage ; les autres réglages ne se départagent qu'à quelques centièmes.

\reportfig{fig_lightgbm_recherche}{fig:lgb}{Recherche aléatoire de 60 configurations de LightGBM
(validation sur 21\,000 lignes). Chaque point est une configuration, le trait noir la médiane par valeur,
les tirets orange la configuration de départ.}

\textbf{CatBoost}~[5] utilise des arbres symétriques (tous les nœuds d'une même profondeur posent la même
question) et code chaque catégorie par une moyenne lissée de la cible, calculée sur les seules lignes qui précèdent
dans une permutation aléatoire. Parmi les configurations comparées sur la validation, la meilleure
(profondeur 7, $\eta = 0{,}06$, $\lambda = 8$, $p = 1$, 3\,150 arbres) inclut une variable de contexte
fédération\,|\,pays\,|\,région\,|\,année ; sans elle, l'erreur de validation passe de 78,18 à 78,30\kg.

\subsection{TabM : un ensemble de réseaux dans un seul réseau}
TabM~[6] entraîne simultanément $k$ perceptrons qui partagent l'essentiel de leurs poids et moyenne leurs
prédictions. Chaque variable numérique est d'abord codée par morceaux~[7], ce qui permet au réseau de traiter
différemment chaque plage de valeurs. Nous avons tiré 16 configurations (2 à
4 couches de 256 à 512 neurones, $k = 16$ ou 32, taux de $5\cdot10^{-4}$ à $2\cdot10^{-3}$, dropout, nombre
d'intervalles, pénalité des poids, exposant $p$), plus notre configuration de départ (3 couches de
384 neurones, $k = 32$, 64 intervalles), chacune arrêtée après 20 époques sans progrès en validation. Les
17 erreurs vont de 77,77 à 78,23\kg{} et la meilleure est la configuration de départ. Ces écarts sont de
l'ordre de la variabilité d'entraînement : la même configuration, réentraînée avec 38 époques au lieu de
40 et sur une autre carte graphique, donne 78,24\kg{} en validation croisée au lieu de 78,12, et sur le test
interne deux entraînements donnent 78,16 et 79,04\kg. Pour TabM, moyenner plusieurs entraînements serait un
meilleur usage du calcul qu'une recherche plus large.

\subsection{TabPFN : un modèle préentraîné qui apprend dans le contexte}
TabPFN~[8,9] est un transformer préentraîné une fois pour toutes, par ses auteurs, sur un très grand nombre de
jeux de données simulés, à approcher la prédiction bayésienne de la cible de lignes à partir d'autres lignes
étiquetées. Pour prédire, on lui donne nos lignes étiquetées \emph{en entrée}, avec les lignes à prédire ; ses
paramètres ne changent pas. Il renvoie une loi pour $y$, dont nous prenons l'espérance. Nous utilisons la version 3.5 du modèle (bibliothèque \code{tabpfn} 9.1.0), avec les 16 variables
d'origine et tout le contexte disponible (67\,200 lignes en validation croisée, 105\,000 pour le test).
Le modèle est appliqué plusieurs fois à des présentations différentes des données (ordre des colonnes,
prétraitements), puis on moyenne ; nous avons fait varier ce nombre de passages (paramètre \code{n\_estimators}) et la cible sur le pli 1.
Passer de 1 à 4 passages réduit l'erreur de 0,30\kg, passer de 4 à 16 de 0,07, pour un temps de calcul
proportionnel au nombre de passages. À nombre de passages égal, la cible ramenée au poids gagne 0,07 à
0,09\kg. Nous retenons 8 passages et la cible ramenée ($p = 2/3$) : 77,08\kg{} en validation croisée avec
4 passages et la cible en kilogrammes, 76,94 avec ce choix.

\subsection{Comparaison}
\begin{tableau}\small
\begin{tabular}{@{}lccc@{}}
\toprule
Modèle & Variables & Validation croisée & Test interne \\
\midrule
Ridge ($p = 1$) & 16 & 83,93 & 83,81 \\
CatBoost & 36 & 78,40 & 78,32 \\
LightGBM, configuration de départ & 32 & 78,23 & 78,10 \\
LightGBM, meilleure configuration ($\eta = 0{,}02$) & 32 & 78,10 & 77,98 \\
LightGBM, configuration de départ, 5 graines & 32 & 78,07 & 77,89 \\
\code{HistGradientBoosting}, 5 graines & 54 & 78,26 & 78,01 \\
TabM & 24 & 78,12 & 78,16 \\
\textbf{TabPFN (8 passages, cible ramenée)} & 16 & \textbf{76,94} & \textbf{76,78} \\
\bottomrule
\end{tabular}\par\vspace{2pt}
\legende{tab:individuels}{RMSE (kg) des modèles individuels. Validation croisée : mêmes cinq plis sur
84\,000 lignes. Test interne : apprentissage sur les 84\,000 lignes, mesure sur les 21\,000 lignes mises de côté.}
\end{tableau}

La figure~\ref{fig:courbes} prédit les 16\,800 lignes du pli 1 en n'apprenant que sur une partie des
67\,200 autres. Entre 30\,000 et 67\,200 lignes, Ridge ne gagne plus que 0,6\kg{} ; les trois autres modèles
gagnent encore 1,8 à 2,0\kg{} et leurs courbes ne s'aplatissent pas, d'où l'intérêt de réentraîner les modèles
finaux sur les 105\,000 lignes. LightGBM et TabM sont indiscernables à toutes les tailles. TabPFN est devant
partout, mais son avance fond quand les données augmentent : 3,6 à 3,8\kg{} avec 1\,000 lignes, 1,4\kg{} avec
67\,200 ; avec 1\,000 lignes, il fait aussi bien que Ridge avec 67\,200. Son avantage est de même nettement plus grand sur
les profils rares (tableau~\ref{tab:rarete}). Nous l'interprétons comme l'effet de l'a priori appris au
préentraînement, qui compte d'autant plus que les données sont rares.

\reportfig[0.68\linewidth]{fig_courbes_apprentissage}{fig:courbes}{Erreur sur les 16\,800 lignes du pli 1
selon le nombre de lignes d'apprentissage (sous-échantillons emboîtés).}

\begin{tableau}\small
\begin{tabular}{@{}lccc@{}}
\toprule
Écart de RMSE avec TabPFN (kg) & LightGBM & CatBoost & TabM \\
\midrule
Profils rares : groupes de 1 à 4 lignes (8\,600 lignes) & +2,2 [1,6 ; 2,8] & +2,9 [2,3 ; 3,5] & +2,0 [1,3 ; 2,5] \\
Profils fréquents : groupes de 500 lignes et plus & +0,6 [0,2 ; 1,0] & +0,2 [$-0{,}1$ ; 0,6] & +0,7 [0,4 ; 1,0] \\
\bottomrule
\end{tabular}\par\vspace{2pt}
\legende{tab:rarete}{RMSE de chaque modèle moins RMSE de TabPFN (4 passages, cible en kg), avec intervalles à
95\,\% ; un écart positif signifie que TabPFN fait mieux. Le groupe d'une ligne est sa combinaison
sexe\,|\,équipement\,|\,catégorie\,|\,fédération, et sa taille est comptée sur les 84\,000 lignes.}
\end{tableau}

% =====================================================================
\section{Où et pourquoi les modèles se trompent}\label{sec:erreurs}
Les erreurs des six modèles de notre premier mélange ont des corrélations de 0,96 à 0,99 ; leurs
prédictions ne s'écartent en moyenne que de 11\kg{} les unes des autres, pour une erreur absolue moyenne de
58\kg. Les 5\,\% de lignes les plus mal prédites portent 35\,\% de la somme des carrés des erreurs. Sur
les lignes où l'erreur du mélange dépasse deux fois sa RMSE (5\,\% des lignes), les six modèles se trompent
tous dans le même sens : tous surestiment le total, ou tous le sous-estiment (83\,\% des lignes dans
l'ensemble). Sur ces lignes, l'erreur de tout mélange à poids positifs de somme 1 est au moins égale, en
valeur absolue, à la plus petite des six : combiner les modèles ne peut pas corriger ces grosses erreurs. Notons $X$ les variables d'une ligne et $f^*(X)=\mathbb{E}[Y\mid X]$. Pour toute prédiction
$f(X)$ calculée à partir de ces seules variables, le terme croisé étant nul ($\mathbb{E}[Y-f^*(X)\mid X]=0$),
\[\mathbb{E}\big[(Y-f(X))^2\big]=\mathbb{E}\big[(f^*(X)-f(X))^2\big]+\mathbb{E}\big[\operatorname{Var}(Y\mid X)\big].\]
Le premier terme dépend du modèle ; le second, la dispersion des totaux entre lignes aux variables identiques,
n'en dépend pas. Il est grand : trois hommes aux variables presque identiques (même fédération autrichienne,
\emph{Raw}, catégorie 120\kg, même année, 27 à 30 ans, 107 à 113\kg{} de poids de corps) ont réalisé 777,5, 617,5 et 237,5\kg ; les six modèles
leur prédisent tous entre 612 et 660\kg. Les prédictions sont par ailleurs bien calibrées (pente de la
régression du total sur la prédiction entre 0,99 et 1,01 pour tous les modèles). Ce qui manque n'est donc pas
un meilleur algorithme mais une information absente des variables de la ligne : le niveau propre de
l'athlète. La borne $\mathbb{E}[\operatorname{Var}(Y\mid X)]$ ne vaut que pour une prédiction fondée sur les
seules variables de la ligne ; la section~\ref{sec:appariement} la contourne en utilisant les totaux des autres
lignes qui proviennent probablement du même athlète.

% =====================================================================
\section{Combiner les modèles}\label{sec:melange}
Les poids d'un mélange des prédictions hors pli $\hat y_{ij}$ résolvent
$\min_w \sum_i \big(y_i - \sum_j w_j \hat y_{ij}\big)^2$ sous $w_j \ge 0$, $\sum_j w_j = 1$ : c'est le stacking de
Wolpert~[15] dans la forme de Breiman~[11], résolu par SLSQP~[12]. Avant le mélange, chaque prédiction est
rapprochée du total moyen des lignes d'apprentissage \emph{identiques} sur les 16 variables, quand il en
existe : $\hat y \leftarrow 0{,}25\,\hat y + 0{,}75\,\bar y_{\text{profil}}$. L'erreur quadratique d'un tel mélange
est la moyenne pondérée des erreurs des modèles, diminuée de leur dispersion autour du mélange~[10]. Pour
notre premier mélange de six modèles, en moyenne sur les lignes, cela donne $6\,015{,}3 - 124{,}9 = 5\,890{,}4$,
soit 76,75\kg{} de RMSE. Des modèles aussi corrélés ne peuvent se compenser
que marginalement ; pour deux modèles, le mélange ne bat le meilleur que si la corrélation de leurs erreurs
est inférieure au rapport de leurs RMSE, la plus petite sur la plus grande~[13], et entre TabPFN et chaque autre modèle elle l'est de peu
(0,967 à 0,972 contre 0,983 à 0,988).

Évalué en validation imbriquée, ce mélange fait 76,76\kg ; trois changements le portent à 76,65\kg{} : la nouvelle version de
TabPFN remplace l'ancienne ; les deux LightGBM sont moyennés sur 5 graines et un \code{HistGradientBoosting}
entraîné sur 54 variables, dont des moyennes de la cible par cohorte, ajoute des erreurs un peu moins
corrélées ; les poids diffèrent selon que le groupe de la ligne est rare ou fréquent (TabPFN reçoit 68\,\% du
poids sur les profils rares, 51\,\% sur les fréquents), ce qu'appelle le tableau~\ref{tab:rarete}. Le mélange
gagne ainsi 0,28\kg{} sur TabPFN seul, et dans chacun des cinq plis (de 0,20 à 0,40\kg{} selon le pli) : un gain réel, mais modeste.

% =====================================================================
\section{Correction par appariement}\label{sec:appariement}

\subsection{Principe}
Aucune variable n'identifie l'athlète, mais une ligne donne son sexe, son pays, sa région et son année de
naissance, déduite de l'âge et de l'année de la compétition à un an près. Les lignes étiquetées qui partagent
ces attributs, avec le même équipement et un poids de corps proche, ont de bonnes chances d'être
\emph{celles du même athlète} à d'autres compétitions. C'est un problème d'appariement d'enregistrements
(\emph{record linkage}~[17]) : nous ne savons pas quelles lignes concernent la même personne, mais nous
pouvons apprendre à quel point il est probable qu'elles le fassent.

Pour une ligne $i$ de prédiction $\hat y_i$, nous cherchons dans un réservoir de lignes étiquetées les
\emph{candidates} $c$ de même sexe, pays, région et année de naissance compatible, puis appliquons six filtres
de plus en plus stricts : aucun ; même équipement et poids à 10\,\% près ; poids à 5\,\% près et compétition à
3 ans près ; poids à 3\,\% près et même fédération ; poids à 2\,\% près et compétition à un an près ; équipement
différent. Pour chaque filtre, avec $n$ candidates, nous calculons
\[
d_c = y_c \Big(\frac{b_i}{b_c}\Big)^{2/3} - \hat y_i, \qquad
m_\kappa = \frac{n}{n + \kappa}\cdot\frac1n \sum_c d_c \quad (\kappa \in \{0{,}5 ; 2 ; 8\}),
\]
l'écart moyen, rétréci vers zéro quand les candidates sont peu nombreuses, entre leur total ramené au poids de
la ligne et la prédiction ; s'y ajoutent $\log(1+n)$, les écarts moyens de poids et d'année et la dispersion des
$d_c$. Un LightGBM de second niveau (15 feuilles, 200 arbres, $\eta = 0{,}03$, au moins 300 lignes par feuille,
moyenne de 5 graines) apprend l'erreur $y_i - \hat y_i$ à partir de ces 43 caractéristiques. Seuls les totaux
réels des candidates et la prédiction de la ligne elle-même y entrent, comme pour une ligne de test. En
validation croisée, pour évaluer un pli, le modèle de second niveau est entraîné sur les autres plis avec des
caractéristiques calculées sans aucune ligne de ce pli, puis appliqué au pli avec pour réservoir les quatre
autres. Pour le test, il est entraîné sur les 84\,000 lignes de développement, qui forment le réservoir.

\subsection{Résultats}
\begin{tableau}\small
\begin{tabular}{@{}lccc@{}}
\toprule
Prédiction corrigée & Avant & Après & Gain [IC 95\,\%] \\
\midrule
TabPFN seul & 76,91 & 76,55 & $-0{,}36$ [$-0{,}42$ ; $-0{,}31$] \\
Mélange avec l'ancienne version de TabPFN & 76,69 & 76,29 & $-0{,}39$ [$-0{,}45$ ; $-0{,}34$] \\
\textbf{Mélange retenu} & \textbf{76,64} & \textbf{76,27} & $\mathbf{-0{,}37}$ [$-0{,}43$ ; $-0{,}32$] \\
\bottomrule
\end{tabular}\par\vspace{2pt}
\legende{tab:appariement}{RMSE (kg) en validation croisée sur 84\,000 lignes ; « Après » est évalué en validation
imbriquée. Le gain est positif dans chacun des cinq plis.}
\end{tableau}

Le gain est presque le même quelle que soit la prédiction corrigée : la correction n'exploite pas la diversité
des modèles, elle apporte une information qu'aucun n'avait. Plusieurs analyses indiquent qu'il s'agit bien de
l'historique de l'athlète. Avec les seules candidates de même génération et de même région, sans condition
d'équipement ni de poids, le gain est nul ($-0{,}03$ [$-0{,}07$ ; 0,01]) ; avec les seules candidates proches
(même équipement, poids proche), il atteint $-0{,}32$. Le gain croît avec l'âge : 0,07\kg{} avant 24 ans, 1,50\kg{}
après 50 ans (figure~\ref{fig:appariement}). Exemple : un homme de 68 ans, 92,6\kg, \emph{Single-ply}, fédération
anglaise EPA, a réalisé 672,5\kg{} en 2019 ; le mélange lui prédisait 518\kg. Trois candidates (58, 60 et 64 ans,
91,9 à 92,7\kg, \emph{Single-ply}, fédérations EPA et BP, 2009 à 2016) ont réalisé 690 à 712,5\kg, et la correction
porte la prédiction à 582\kg. Elle se trompe aussi : un homme de 54 ans et 122,9\kg{} a réalisé 645\kg{} quand ses
deux candidates les plus proches avaient fait 872,5 et 907,5\kg, et sa prédiction passe de 736 à 810\kg.

\textbf{Quand faire confiance aux candidates.} Une candidate n'est pas forcément le même athlète. La
figure~\ref{fig:appariement} donne, par filtre et par taille de cohorte (nombre de lignes de même sexe, pays,
région et année de naissance), la pente de la régression de l'erreur du mélange sur l'écart moyen des
candidates à la prédiction : 1 si cet écart se retrouvait en moyenne tout entier dans l'erreur, 0 si les
candidates n'apportaient rien. La pente vaut 0,47 pour des candidates à 2\,\% du poids et à un an près dans une
cohorte de 1 à 3 personnes, et 0,00 à 0,03 dans les cohortes de plus de 50 personnes, où les candidates sont
presque toujours d'autres athlètes. C'est la logique de l'appariement probabiliste~[17] : un accord sur des
attributs rares pèse plus qu'un accord sur des attributs fréquents. Le modèle de second niveau apprend cette
pondération ; sur le mélange avec l'ancienne version de TabPFN, une régression linéaire sur les mêmes
caractéristiques ne gagne que 0,29\kg, contre 0,37 à 0,39 pour LightGBM, XGBoost, CatBoost ou
\code{HistGradientBoosting}. Enfin, donner directement à LightGBM les
totaux des candidates, leur nombre et leur proximité, en plus de ses 32 variables, ne réduit son erreur que de
0,08\kg{} (intervalle $[-0{,}22 ; 0{,}05]$) : l'information est mieux exploitée comme correction d'une
prédiction que comme variable d'entrée.

\reportfig{fig_appariement}{fig:appariement}{À gauche, pente de l'erreur du mélange sur l'écart moyen des
candidates à la prédiction, par filtre et taille de cohorte (candidates prises dans les autres plis ; « l. » :
nombre de lignes). À droite, baisse de la RMSE apportée par la correction, par âge.}

% =====================================================================
\section{Résultats et soumissions retenues}\label{sec:resultats}

\begin{tableau}\small
\begin{tabular}{@{}lccc@{}}
\toprule
Solution & Validation croisée & Test interne & Public \\
\midrule
TabPFN seul & 76,94 & 76,78 & 76,42 \\
TabPFN + correction par appariement & 76,55 & 76,08 & 75,86 \\
\textbf{Mélange + correction par appariement (solution finale)} & \textbf{76,27} & \textbf{75,83} & \textbf{75,65} \\
\bottomrule
\end{tabular}\par\vspace{2pt}
\legende{tab:soumissions}{RMSE (kg). Validation croisée sur 84\,000 lignes (imbriquée pour la correction) ; test
interne : tout est ajusté sur les 84\,000 lignes et mesuré une fois sur les 21\,000 autres ; public : environ
13\,500 lignes du test Kaggle.}
\end{tableau}

Sur le test interne, la solution finale fait 0,95\kg{} de moins que TabPFN seul (intervalle $[0{,}79 ; 1{,}11]$) :
le mélange en apporte 0,27 (76,50\kg{} avant correction) et la correction 0,68 ($[0{,}55 ; 0{,}80]$). La
correction y rapporte presque deux fois plus qu'en validation croisée, où le modèle de second niveau apprend
sur 67\,200 lignes et cherche ses candidates parmi 67\,200 lignes, contre 84\,000 ici. En ne gardant comme
réservoir qu'une partie des 84\,000 lignes, tirée au hasard, le gain sur le test interne vaut 0,23\kg{} avec
33\,600 lignes, 0,45 avec 50\,400, 0,59 avec 67\,200 et 0,68 avec 84\,000 : plus on connaît de participations,
plus on a de chances d'y retrouver l'athlète. Nous avons donc soumis une variante dont le réservoir contient les
105\,000 lignes étiquetées. Sur le classement public, elle est à égalité avec la solution finale (75,650 contre
75,651\kg) : si l'effet existe sur le test Kaggle, il est trop petit pour être vu sur 13\,500 lignes. Nous
retenons ces deux soumissions pour le classement final.

% =====================================================================
\section{Conclusion}
Sur ce problème, l'information compte plus que l'algorithme. Les modèles individuels se distinguent surtout
sur les profils rares, où le préentraînement de TabPFN fait la différence ; les réglages de LightGBM et de TabM
ne se départagent qu'à quelques centièmes ou dans le bruit d'entraînement, et tous les modèles butent sur la
même limite, le niveau propre de l'athlète. Les combiner rapporte peu, ce que leurs corrélations expliquent.
Retrouver dans les données les participations probables du même athlète rapporte plus que le mélange
(0,37 contre 0,28\kg{} en validation croisée, 0,68 contre 0,27\kg{} sur le test interne).

\textbf{Limites.} La correction repose sur des identités probables : elle se trompe quand deux athlètes
partagent les mêmes attributs, ou quand un athlète réalise une performance inhabituelle. Les poids du mélange
et les réglages de la correction ont été choisis sur les 84\,000 lignes de développement ; le test interne en
donne une mesure indépendante. Enfin, sans identifiant d'athlète, nous ne savons pas mesurer la précision sur
des athlètes entièrement nouveaux.

% =====================================================================
{\small
\section*{Références}
\begin{enumerate}[label={[\arabic*]},leftmargin=2.2em]
\item A.\,E.~Hoerl, R.\,W.~Kennard. Ridge regression: biased estimation for nonorthogonal problems. \emph{Technometrics} 12(1), 1970.
\item J.\,H.~Friedman. Greedy function approximation: a gradient boosting machine. \emph{Annals of Statistics} 29(5), 2001.
\item T.~Chen, C.~Guestrin. XGBoost: a scalable tree boosting system. \emph{KDD}, 2016.
\item G.~Ke et al. LightGBM: a highly efficient gradient boosting decision tree. \emph{NeurIPS}, 2017.
\item L.~Prokhorenkova et al. CatBoost: unbiased boosting with categorical features. \emph{NeurIPS}, 2018.
\item Y.~Gorishniy, A.~Kotelnikov, A.~Babenko. TabM: advancing tabular deep learning with parameter-efficient ensembling. \emph{ICLR}, 2025.
\item Y.~Gorishniy, I.~Rubachev, A.~Babenko. On embeddings for numerical features in tabular deep learning. \emph{NeurIPS}, 2022.
\item S.~Müller, N.~Hollmann, S.~Pineda Arango, J.~Grabocka, F.~Hutter. Transformers can do Bayesian inference. \emph{ICLR}, 2022.
\item N.~Hollmann et al. Accurate predictions on small data with a tabular foundation model. \emph{Nature} 637, 2025.
\item A.~Krogh, J.~Vedelsby. Neural network ensembles, cross validation, and active learning. \emph{NIPS 7}, 1995.
\item L.~Breiman. Stacked regressions. \emph{Machine Learning} 24(1), 1996.
\item D.~Kraft. A software package for sequential quadratic programming. DFVLR-FB 88-28, 1988.
\item J.\,M.~Bates, C.\,W.\,J.~Granger. The combination of forecasts. \emph{Operational Research Quarterly} 20(4), 1969.
\item G.~Marković, D.~Sekulić. Modeling the influence of body size on weightlifting and powerlifting performance. \emph{Collegium Antropologicum} 30(3), 2006.
\item D.\,H.~Wolpert. Stacked generalization. \emph{Neural Networks} 5(2), 1992.
\item OpenPowerlifting, bibliothèque \code{opltypes} 0.2.0, \code{src/age.rs}. \url{https://docs.rs/crate/opltypes/0.2.0/source/src/age.rs}
\item I.\,P.~Fellegi, A.\,B.~Sunter. A theory for record linkage. \emph{Journal of the American Statistical Association} 64(328), 1969.
\end{enumerate}
\par}

\end{document}
